\section{F-SM-34: the scalar adjoints do not force free-scalar normalization}

\subsection{Exact conclusion}
There is a complete and cocomplete \(k\)-linear abelian
category \(\mathcal N\), with enough projectives, and an exact
functor \(\kappa:\operatorname{Vect}_k\to\mathcal N\) having
both adjoints, such that the canonical counit
\[
              L\varphi(\kappa k)\longrightarrow k[0]
\]
is an equivalence, whereas the same counit is not an equivalence
on the countably generated free space \(k^{(\mathbb N)}\).
Here \(\varphi\dashv\kappa\), and even the underived
counit \(\varphi\kappa V\to V\) is an isomorphism for every
\(V\).  Countable coproducts of monomorphisms fail to be
monomorphisms in an explicitly exhibited family.

Passing to \(\operatorname{Ch}_{\ge0}(\mathcal N)\) supplies
an exact scalar functor with both adjoints, the same canonical
normalization on \(k\), and a natural monomorphism of every
object into a fully derived-acyclic object.  Even its ordinary
localization conclusion fails.  This is a counterexample to
the \emph{additive, acyclic-effacement version} without a
coproduct condition.  No symmetric monoidal tensor on this
category satisfying the original archive assumptions is
constructed here.  Consequently this result does not refute
the full tensor-and-interval assertion in cote~134-8.

\subsection{Complete construction and proof}

\paragraph{1. The coefficient functor.}
Fix a field \(k\), and set
\[
                       E(V)=V^{**}/V,
\]
where \(V\to V^{**}\) is the canonical evaluation embedding.
The functor \(E\) is \(k\)-linear, additive and exact:
every short exact sequence of vector spaces splits, and an
additive functor preserves a split exact sequence.
It vanishes on finite-dimensional spaces, in particular on
\(k\).

For \(V=k^{(\mathbb N)}\), it is nonzero.  Identify
\(V^*=k^{\mathbb N}\).  The quotient
\(k^{\mathbb N}/k^{(\mathbb N)}\) is nonzero, as witnessed
by the all-ones sequence.  Choose a nonzero linear functional
on that quotient and pull it back to \(k^{\mathbb N}\).
It is a nonzero element of \(V^{**}\) vanishing on the
finite-support subspace.  No nonzero evaluation functional
from \(V\) vanishes on that subspace, since its coordinate
values recover the vector.  Its class in \(E(V)\) is therefore
nonzero.  This uses ordinary vector-space choice only.

\paragraph{2. The two-sort abelian category.}
An object of \(\mathcal N\) is
\[
 X=(X_0,X_1,b:X_0\to X_1,a:E(X_1)\to X_0),
                 \qquad ba=0,\quad aE(b)=0.                  \tag{SM34.1}
\]
A morphism is a pair \((u_0,u_1)\) satisfying
\(u_1b=b'u_0\) and \(u_0a=a'E(u_1)\).
Kernels and cokernels are computed on the two vector spaces:
exactness of \(E\) gives the induced structure maps and
preserves both relations.  The comparison from coimage to
image is an isomorphism componentwise.  Thus \(\mathcal N\)
is abelian, and the forgetful functor to
\(\operatorname{Vect}_k^2\) is exact and faithful.

For vector spaces \(U,V\), define
\[
 \mathsf P(U,V)=
       (U\oplus E(V),\ U\oplus V),
 \quad b(u,e)=(u,0),\quad a(e_U,e_V)=(0,e_V).                \tag{SM34.2}
\]
Here \(E(U\oplus V)=E(U)\oplus E(V)\).  The two relations
are immediate.  A map \(\mathsf P(U,V)\to X\) is uniquely
specified by maps \(u:U\to X_0\), \(v:V\to X_1\):
its components are
\((u,aE(v))\) and \((bu,v)\).  Hence
\[
 \operatorname{Hom}_{\mathcal N}(\mathsf P(U,V),X)
     =\operatorname{Hom}_k(U,X_0)\times\operatorname{Hom}_k(V,X_1).
                                                               \tag{SM34.3}
\]
This is a free-object adjunction.  Since the forgetful functor
is exact, each \(\mathsf P(U,V)\) is projective.  The counit
\(\mathsf P(X_0,X_1)\to X\) is componentwise epic: its
components contain the identity of \(X_0\) and of \(X_1\),
respectively.  There are therefore enough projectives.

For completeness, all small limits and colimits exist.
Products are formed componentwise; the map
\(E(\prod X_{i,1})\to\prod X_{i,0}\) has components
\(a_iE(\operatorname{pr}_i)\).  Both relations follow
after projection to each factor.  Together with the already
constructed kernels this gives all limits.

Here is an explicit construction of coproducts, without
assuming that \(E\) preserves them.  Write \(U\) for the
forgetful functor, \(\mathsf P\) for its left adjoint, and
\(T=U\mathsf P\).  For a family \(X_i\), put
\(L=\bigoplus_iUX_i\), \(R=\bigoplus_iTUX_i\) in
\(\operatorname{Vect}_k^2\).  Form the coequalizer in
\(\mathcal N\)
\[
             \mathsf P(R)\rightrightarrows\mathsf P(L)
                                      \longrightarrow Q.     \tag{SM34.4}
\]
The first arrow is induced by the sum of the algebra action
maps \(TUX_i\to UX_i\).  The second is the free extension
of the map \(R\to T L\) whose restriction to the \(i\)-th
summand is \(T\) of the summand inclusion \(UX_i\to L\).
These action maps are simply the free-adjunction counits.
One may also verify them by the square-zero monad
\(T=\mathrm{id}\oplus H\), \(H(U,V)=(E(V),U)\): its
multiplication adds the two \(H\)-summands and kills
\(H^2\).  Associativity adds the three linear \(H\)-terms
and kills all higher terms on either side; its algebra
relation is exactly (SM34.1).

Coequalizers exist as cokernels of differences.  A map from
\(\mathsf P(L)\) to an object \(Z\) corresponds to a family
of underlying linear maps \(UX_i\to UZ\).  Equalizing the
two arrows in (SM34.4) says precisely that each map commutes
with the two structure maps of \(X_i\) and \(Z\).
Thus \(\operatorname{Hom}(Q,Z)=\prod_i
\operatorname{Hom}(X_i,Z)\), proving the coproduct property.
Coproducts and cokernels give all small colimits.

\paragraph{3. Both scalar adjoints.}
Define \(\kappa V=(V,0,0,0)\).  Directly from the
morphism equations,
\[
 \begin{split}
 \operatorname{Hom}_{\mathcal N}(X,\kappa V)
       &=\operatorname{Hom}_k(\operatorname{coker}a,V),\\
 \operatorname{Hom}_{\mathcal N}(\kappa V,X)
       &=\operatorname{Hom}_k(V,\ker b).
 \end{split}                                                 \tag{SM34.5}
\]
Thus \(\varphi X=\operatorname{coker}a\) is left adjoint
to \(\kappa\), and \(\psi X=\ker b\) is its right
adjoint.  The functor \(\kappa\) is exact by componentwise
exactness, preserves all coproducts and products by these
adjoints, and is the vector-space copower of its distinguished
object \(\kappa k\).  The underived counit is literally
\(\varphi\kappa V=V\).

Put
\[
 P_0(V)=\mathsf P(V,0)=(V,V,\mathrm{id},0),\qquad
 P_1(V)=\mathsf P(0,V)=(E(V),V,0,\mathrm{id}).                \tag{SM34.6}
\]
The epimorphism \(P_0(V)\to\kappa V\) has kernel
\(S(V)=(0,V,0,0)\).  The epimorphism
\(P_1(V)\to S(V)\) has kernel \(\kappa E(V)\).
Repeating gives a projective resolution of \(\kappa V\)
with terms
\[
       P_{2n}=P_0(E^nV),\qquad P_{2n+1}=P_1(E^nV).
                                                               \tag{SM34.7}
\]
The odd differential has components \((0,\mathrm{id})\);
the positive even differential has components
\((\mathrm{id},0)\).  The preceding kernel calculations
prove exactness at every term, not merely that the
differentials compose to zero.

Applying \(\varphi\) gives \(E^nV\) in degree \(2n\)
and zero in odd degrees.  Therefore
\[
 L_{2n}\varphi(\kappa V)=E^nV,\qquad
 L_{2n+1}\varphi(\kappa V)=0\quad(n\ge0).                    \tag{SM34.8}
\]
For \(V=k\), this is \(k\) in degree zero and zero
elsewhere, and its canonical counit is the identity in
degree zero.  For \(V=k^{(\mathbb N)}\), degree two is
the nonzero space \(E(V)\).  This proves failure of the
free-scalar counit despite both scalar adjoints and the
canonical unit normalization.

There is also an explicit failure of exact countable sums.
Since \(E(k)=0\), the maps
\[
                S(k)=P_1(k)\hookrightarrow P_0(k)
\]
with components \((0,\mathrm{id})\) are monomorphisms.
Their countable coproduct is
\[
                    P_1(V)\longrightarrow P_0(V),
                        \qquad V=k^{(\mathbb N)},            \tag{SM34.9}
\]
because the free functor \(\mathsf P\) preserves
coproducts.  Its components are again
\((0,\mathrm{id})\), by the free-map formula (SM34.3).
Its kernel is \(\kappa E(V)\ne0\), so it is not monic.
Likewise \(\bigoplus_nL\varphi(\kappa k)\to
L\varphi(\bigoplus_n\kappa k)\) fails in degree two.

\paragraph{4. Actual acyclic effacement after adjoining chain degrees.}
Let \(\mathcal M=\operatorname{Ch}_{\ge0}(\mathcal N)\).
It is complete and cocomplete abelian.  It has enough
projectives: for projective \(Q\in\mathcal N\), the
complex \(Q[0]\) is projective in the chain category, as
is each disk \(D^n(Q)\), \(n\ge1\); their Hom functors
are evaluation in degrees \(0\) and \(n\), respectively,
followed by \(\operatorname{Hom}(Q,-)\).  Choose projective
epimorphisms onto each term of a complex and use the
corresponding disks and degree-zero complex.  Their
coproduct maps degreewise epimorphically to the complex.
There are at most two disk summands in any fixed degree
in this chosen epimorphism, so its construction does not
assume exact infinite coproducts.  A projective chain
object is a retract of such a coproduct.  The disk part
is explicitly contractible, and the remaining degree-zero
part is projective.

The exact scalar functor is
\(\kappa_{\mathcal M}V=(\kappa V)[0]\).  Its left and
right adjoints are, respectively,
\[
       \varphi_{\mathcal M}(C)=\varphi(H_0C),\qquad
                     \psi_{\mathcal M}(C)=\psi(C_0).           \tag{SM34.10}
\]
The adjunctions follow by factoring chain maps to a
degree-zero complex through \(H_0\), and by evaluating
chain maps from a degree-zero complex in degree zero.

There is a canonical identification
\[
       L\varphi_{\mathcal M}(C)=\mathbb L\varphi(C)
                                  \quad\text{in }D_{\ge0}(k), \tag{SM34.11}
\]
where the right-hand side regards \(C\) as a complex in
\(\mathcal N\).  To verify it, resolve \(C\), as an
object of the chain category, by projective chain objects
\(R_{r,\bullet}\).  Each has projective terms,
no positive chain homology, and projective \(H_0\), since
it is a retract of a coproduct of the disks and degree-zero
projectives just described.  The first-quadrant total
complex \(\operatorname{Tot}R\to C\) is a quasi-isomorphism
by exactness in the resolution direction.  The row
augmentations also give a quasi-isomorphism
\(\operatorname{Tot}R\to H_0(R_{\bullet,\bullet})\).
Both assertions use finite diagonals in each total degree.
The latter complex has projective \(\mathcal N\)-terms,
and applying \(\varphi\) gives exactly the resolution
formula for \(L\varphi_{\mathcal M}(C)\).  This proves
(SM34.11), with its canonical comparisons.  No infinite
sum of short exact sequences in \(\mathcal N\) is used.

Define \(A(C)_n=C_n\oplus C_{n-1}\), with \(C_{-1}=0\),
and
\[
        d(x,y)=(dx+y,-dy),\quad
        a_C(x)=(x,0),\quad h(x,y)=(0,x).                     \tag{SM34.12}
\]
These finite-matrix formulas make sense in any abelian
category.  The map \(a_C:C\to A(C)\) is a natural
degreewise monomorphism, and \(dh+hd=1\).
Thus \(A(C)\) is contractible and (SM34.11) implies
\(L\varphi_{\mathcal M}(A(C))=0\), including degree zero.
The canonical scalar counit on \(k\) still is an
equivalence by (SM34.8), and the one on
\(k^{(\mathbb N)}\) still is not.

\paragraph{5. Why the additive localization conclusion actually fails.}
Put
\[
 W_{\mathcal M}=(L\varphi_{\mathcal M})^{-1}
                         (\text{isomorphisms}).
\]
Suppose, contrary to the claim,
that the induced ordinary functor
\(W_{\mathcal M}^{-1}\mathcal M\to D_{\ge0}(k)\)
were an equivalence.  Since quasi-isomorphisms belong to
\(W_{\mathcal M}\), localization in stages and (SM34.11)
would make
\[
          \mathbb L\varphi:D_{\ge0}(\mathcal N)
                                         \longrightarrow D_{\ge0}(k)
\]
a localization functor.  It has right adjoint the exact
termwise functor \(\kappa\), since
\(\varphi\dashv\kappa\) and \(\varphi\) takes
projectives to projectives.

A localization functor having a right adjoint has fully
faithful right adjoint.  One can check this elementary fact
directly: for \(q:\mathcal A\to\mathcal B\),
\(q\dashv r\), each \(rY\) is local for the inverted
maps, by the adjunction.  A zigzag into \(rY\) can be
uniquely reduced to an actual map by successively extending
across each inverted map; hence
\(\operatorname{Hom}_{\mathcal A}(X,rY)=
\operatorname{Hom}_{\mathcal B}(qX,qrY)\).
The adjunction identifies composition with the counit
\(qrY\to Y\) with a bijection on these Hom sets.
Essential surjectivity of \(q\) and Yoneda make the
counit an isomorphism.  This is also the ordinary instance
of Cisinski Proposition~7.1.17.

In the present example the counit on
\(V=k^{(\mathbb N)}\) has nonzero degree-two source
homology \(E(V)\), by (SM34.8).  It cannot be an
isomorphism in the derived category.  This contradiction
proves failure of the asserted additive localization,
and therefore also of an \(\infty\)-categorical
equivalence induced by this functor. \(\square\)

\subsection{What this counterexample does not certify}
The complete construction certifies the two adjoints of the
scalar functor, its exactness, enough projectives, all limits and colimits,
canonical unit normalization, actual monic acyclic
effacement, and failure of the localization conclusion.
It does not supply the associative, symmetric, unital
tensor of the handwritten original, or identify (SM34.12)
with tensoring by an object of such a monoidal category.
That additional construction remains a substantive gate.
In particular, an exact closed tensor compatible with
the distinguished scalar unit would be impossible here:
F-SM-35 would imply AB4, contradicting the explicit
monomorphism calculation (SM34.9), which also persists in
degree zero in the chain category.

\subsection{Sources}
All example-specific assertions are proved above.  The
derived-category and localization background is the same
as in F-SM-32: Lurie (2017), \emph{Higher algebra},
Theorem~1.3.4.4, p.~109,
\url{https://people.math.harvard.edu/~lurie/papers/HA.pdf};
and Cisinski (2019, author version July~30, 2026),
\emph{Higher categories and homotopical algebra},
Proposition~7.1.17, p.~331,
\url{https://cisinski.app.uni-regensburg.de/CatLR.pdf}.
The scalar/tensor historical distinction refers to
Grothendieck's cote~134-8, footers~46--48,
\url{https://grothendieck.umontpellier.fr/134-8.pdf}.
